\section{A fixed positive angular selection and near-source collisions}
\label{angselect:section}

The following estimate uses angular geometry absent from the scalar
profile. It does not by itself improve the dimension condition: the
far-source estimate specified below remains a separate hypothesis.

Let \(\mu,\nu\) be compactly supported planar probabilities with
\(a_0\le |x-y|\le b_0\) on their product support, where \(a_0>0\).
Assume \(\mu(B(x,r))\le C_F r^s\), with \(s>1\). Write
\(\Theta_x(y)=(y-x)/|y-x|\), and suppose
\begin{equation}\label{angselect:moment}
 (\Theta_x)_*\nu=\rho_x(\theta)\,d\theta,
 \qquad
 B=\int\!\int_{S^1}\rho_x(\theta)^q\,d\theta\,d\mu(x)<\infty,
 \qquad q>1.
\end{equation}
The densities are needed only for \(\mu\)-almost every source point.
They have a jointly measurable version by the Radon--Nikodym theorem
on source space times the circle. If both Frostman exponents exceed
one, the averaged radial-projection theorem supplies
\eqref{angselect:moment} after reducing to a common exponent
\cite[Theorem 3.7]{giow}.

Fix \(A\ge1\), independently of every subsequent resolution, and set
\[
 a_A(x,y)=\mathbf1_{\{\rho_x(\Theta_x(y))\le A\}},
 \qquad \eta_y^A(dx)=a_A(x,y)\,d\mu(x).
\]
Assign selection zero on the exceptional null set of source points.
These are positive measurable kernels, and
\begin{equation}\label{angselect:mass}
 \iint(1-a_A)\,d\mu\,d\nu\le B A^{1-q},
 \qquad
 (\Theta_x)_*(a_A(x,\cdot)\nu)\le A\,d\theta.
\end{equation}
Indeed, the first integral equals
\(\int\!\int_{\{\rho_x>A\}}\rho_x\,d\theta\,d\mu(x)\),
and the second density is exactly
\(\rho_x\mathbf1_{\{\rho_x\le A\}}\).
Thus a sufficiently large fixed \(A\) retains positive total mass.
The loss is fixed; it does not tend to zero as the resolution increases.

\begin{proposition}[Positive near-source collision bound]
\label{angselect:nearprop}
For \(0<\varepsilon<r\le1\), define
\begin{align*}
 \mathcal C^A_{\varepsilon,\le r}
 =\varepsilon^{-1}\iiint
 &\mathbf1_{\{|x-x'|\le r\}}
 \mathbf1_{\{\,||x-y|-|x'-y||\le\varepsilon\}}\\[-2pt]
 &\hspace{8mm}\cdot a_A(x,y)a_A(x',y)
 \,d\mu(x)d\mu(x')d\nu(y).
\end{align*}
Then
\begin{equation}\label{angselect:nearbound}
 \mathcal C^A_{\varepsilon,\le r}
 \le C A\bigl(\varepsilon^{-1}r^{s+1}+r^{s-1}\bigr),
\end{equation}
where \(C\) depends only on the standing geometric and Frostman data.
\end{proposition}

\begin{proof}
Pairs with \(|x-x'|\le\varepsilon\) contribute at most
\(\varepsilon^{-1}\int\mu(B(x,\varepsilon))\,d\mu(x)
\le C_F\varepsilon^{s-1}\).
The diagonal has zero product mass. For
\(u=|x-x'|>\varepsilon\), put
\(e=(x'-x)/u\), \(\ell=|y-x|\), and \(\theta=\Theta_x(y)\).
The squared-distance identity
\[
 |x'-y|^2-|x-y|^2=u^2-2u\ell\,e\cdot\theta
\]
shows that a collision forces
\[
 |e\cdot\theta|\le C_{a_0,b_0}(u+\varepsilon/u).
\]
This set of directions has arclength at most
\(C\min\{1,u+\varepsilon/u\}\). Drop the second selection factor,
which is at most one, and use the first factor's angular domination in
\eqref{angselect:mass}. The retained pin mass for this source pair is
at most \(CA(u+\varepsilon/u)\).

Frostman growth, with dyadic annuli for the second integral, gives
\[
 \int_{0<|x-x'|\le r}|x-x'|\,d\mu(x')\le C_F r^{s+1},
 \qquad
 \int_{0<|x-x'|\le r}\frac{d\mu(x')}{|x-x'|}
 \le C_s C_F r^{s-1}.
\]
Integrate the preceding pin bound and divide by \(\varepsilon\).
The smallest-separation contribution is absorbed because
\(A\ge1\) and \(\varepsilon^{s-1}\le r^{s-1}\).
Both collision factors throughout came from the same positive kernel;
no arbitrary source-pair selection was substituted for that kernel.
\end{proof}

At the parabolic radius the estimate becomes
\begin{equation}\label{angselect:parabolic}
 \mathcal C^A_{\varepsilon,\le\sqrt\varepsilon}
 \le CA\varepsilon^{(s-1)/2}.
\end{equation}
More generally, for fixed
\(1/(s+1)<\alpha<1\) and
\(\kappa=\min\{\alpha(s+1)-1,\alpha(s-1)\}>0\),
\begin{equation}\label{angselect:power}
 \mathcal C^A_{\varepsilon,\le\varepsilon^\alpha}
 \le CA\varepsilon^\kappa.
\end{equation}
These contributions are summable over dyadic resolutions.
One may choose \(\alpha<1/2\), hence control clusters larger than
the parabolic radius. No optimality of this cutoff is asserted.

\begin{proposition}[The exact remaining far-source criterion]
\label{angselect:farcriterion}
Choose \(A\) with \(m_0=1-BA^{1-q}>0\), and \(\alpha\) as above.
Define \(\mathcal C^A_{\varepsilon,>r}\) by replacing
\(|x-x'|\le r\) in Proposition~\ref{angselect:nearprop} by
\(|x-x'|>r\), keeping both weights unchanged.
If some sequence \(\varepsilon_j\to0\) satisfies
\begin{equation}\label{angselect:farhyp}
 \sup_j\mathcal C^A_{\varepsilon_j,>\varepsilon_j^\alpha}<\infty,
\end{equation}
then \(\nu(dy)(d_y)_*\eta_y^A(dt)\) has an
\(L^2(d\nu\,dt)\) density and mass at least \(m_0\).
In particular a positive-\(\nu\) set of pins has positive-length
distance sets from \(\supp\mu\).
\end{proposition}

\begin{proof}
Take a nonnegative smooth mollifier of integral one supported in
\([-1/2,1/2]\). The squared norm of its \(\varepsilon\)-smoothing
of the selected radial law is at most a fixed constant times the full
collision integral of width \(\varepsilon\). Split into its near
and far portions and apply \eqref{angselect:power} and
\eqref{angselect:farhyp}. Weak compactness in the joint \(L^2\)
space and weak convergence of mollifications of the fixed positive
measure identify its density. Equation~\eqref{angselect:mass}
gives the retained mass, and disintegration and Cauchy--Schwarz give
the distance conclusion.
\end{proof}

Only \eqref{angselect:nearbound} is supplied here by angular geometry.
The far hypothesis is not deduced from dimensions. It concerns a fixed
positive pair restriction, unlike the false vanishing-total-loss raw
collision target. Estimating pin mass radially within the angular bands
in the proof is the additional geometric input still required.
